\chapter{Example Episodes and Flips}
\label{app:examples}

This appendix shows real episodes: one for each behavior label, and two
episodes in which an intervention changed the decision. Each was picked by a
fixed rule, stated below, so the examples are illustrations rather than
evidence; the evidence is the counts in Chapter~7. The 20 context paragraphs
are left out. All labels are LLM-derived (Claude Sonnet~4.5), and the labeler
never saw the gold answer (Section~\ref{sec:claude-labeler}).

\section{Label Examples}

These episodes come from the human audit of 100 validation episodes
(Section~\ref{sec:human-audit}). For each label, the example is the audit item
with the lowest ID on which the annotator and Claude both gave that label.
Item IDs were assigned after a seeded shuffle, so this amounts to a random
draw. The last example is the lowest-ID item on which they disagreed.
% Sources: results/human_audit/validation_100/{selection.json,annotations/*.csv,
% scoring_corrected/scores.json,disagreement_review.md};
% results/claude_validation/mixed_final_001/validation/raw_results.jsonl (reasons)

\begin{small}
\begin{longtable}{p{2.3cm}p{12.6cm}}
    \caption{Label examples from the human audit. Each header gives the item
    ID, the source, and which answer the Solver saw as Candidate~A. The human
    label is from the one team member who annotated the item.}
    \label{tab:label-examples} \\
    \toprule
    \endfirsthead
    \toprule
    \endhead
    \bottomrule
    \endfoot
    \multicolumn{2}{l}{\textbf{Adoption} \quad item H014, MuSiQue, A1 listed first} \\
    \midrule
    Question & Who designed the place where Portrait of Francesco delle Opere can be found? \\
    A1 & Giorgio Vasari \\
    C1 & The Uffizi Gallery \\
    A2 & The Uffizi Gallery \\
    Labels & Claude: adoption; human: adoption \\
    Reason & ``A1 states `Giorgio Vasari' as the designer, C1 states `The
        Uffizi Gallery' as the place, and A2 adopts C1's position by stating `The
        Uffizi Gallery'. A1 and C1 represent different positions (designer vs.\
        location), and A2 matches C1's position.'' \\
    \midrule
    \multicolumn{2}{l}{\textbf{Retention} \quad item H005, synthetic, Critic listed first} \\
    \midrule
    Question & In what year did the firm that built the Aldworth Bridge open a
        factory in the city that hosts the Verano Film Festival? \\
    A1 & 1904 \\
    C1 & 2009 \\
    A2 & 1904 \\
    Labels & Claude: retention; human: retention \\
    Reason & ``A1 states `1904' and C1 states `2009', which are different
        positions. A2 states `1904', matching A1's position despite the critic's
        feedback suggesting `2009'. This is retention of the original answer.'' \\
    \midrule
    \multicolumn{2}{l}{\textbf{Agreement} \quad item H003, synthetic, A1 listed first} \\
    \midrule
    Question & Which university did the founder of the oldest newspaper of the city
        housing the Harrowgate Tidal Museum attend? \\
    A1 & Veldhurst University \\
    C1 & Veldhurst University \\
    A2 & Veldhurst University \\
    Labels & Claude: agreement; human: agreement \\
    Reason & ``A1, C1, and A2 all state the same position: `Veldhurst
        University'. Since A1 and C1 already agree, and A2 maintains this same
        position, this is a case of agreement.'' \\
    \midrule
    \multicolumn{2}{l}{\textbf{Out} \quad item H001, MuSiQue, A1 listed first} \\
    \midrule
    Question & In 1971, who introduced the first microprocessor for the production
        company of CyberWorld? \\
    A1 & Intel \\
    C1 & The first microprocessor for the production company of CyberWorld was not
        introduced in 1971. CyberWorld is a 2000 film produced by Intel, and the
        microprocessor referenced in the context of the film is not tied to a 1971
        introduction. Therefore, no company introduced a microprocessor for
        CyberWorld in 1971. \\
    A2 & The evidence is insufficient. \\
    Labels & Claude: out; human: out \\
    Reason & ``A1 provides a specific answer (`Intel'), C1 objects to the
        question's premise (stating no microprocessor was introduced for
        CyberWorld in 1971), and A2 states `The evidence is insufficient' which is
        a different position from both - neither affirming Intel nor making the
        specific objection that C1 makes about the question's premise being
        false.'' \\
    \midrule
    \multicolumn{2}{l}{\textbf{Disagreement} \quad item H002, MuSiQue, Critic listed first} \\
    \midrule
    Question & What genre of music did the composer of Triple Quartet compose? \\
    A1 & Steve Reich \\
    C1 & classical \\
    A2 & classical \\
    Labels & Claude: out; human: adoption \\
    Reason & ``A1 answers `Steve Reich' (a person/composer), C1 answers
        `classical' (a genre), and A2 answers `classical' (a genre). While A2
        matches C1's value, A1 and C1 are not taking different positions on the
        same question - A1 appears to answer `who' while C1 and A2 answer `what
        genre'. Since A1 is not a position on the genre question that C1
        addresses, they cannot be reliably compared as different positions on the
        same matter.'' \\
\end{longtable}
\end{small}

The disagreement shows Claude's most common error in the audit, found in 7 of
the 16 disagreements: it judged whether the two answers were the same
\emph{kind} of answer, which the rules tell it not to do, and called the
episode Out even though A2 repeats C1 word for word. The adoption example has
the same structure, a person against a place, yet Claude labeled it
correctly, so it does not make this error every time. Because it
turns Adoption or Retention into Out rather than into each other, it removes
episodes from the binary target instead of mislabeling them.

\section{Intervention Flips}

These episodes come from the full intervention study (Section~\ref{sec:res-causal}).
The rule, applied in the analysis notebook, takes the four conditions that
survived Holm correction, keeps the episodes in which the targeted edit
flipped the decision and none of the three random edits did, and shows the
one with the shortest text, once for each direction. Such targeted-only flips
occurred in 30 Adoption episodes and 21 Retention episodes. Both examples
below had a first-token probability above 0.999 before the edit, so they are
not typical of where most flips occurred (Section~\ref{sec:res-causal}).
% Sources: results/intervention_layer25/20261106T000808809026Z-02871d677287/merge/labeled_rows.jsonl
% (answers, labels, reasons); question, A1 and C1 from the example cards in
% notebooks/intervention/full_study_results.ipynb, section 1b, which also gives the counts.

\begin{table}[htbp]
    \centering
    \caption{An Adoption-to-Retention flip (synthetic, two-hop). Question:
    ``In what year was the architect of the Marnhill Civic Theatre born?''
    A1: 1927. C1: 1884. All edits have the same size, $0.2\lVert x\rVert$.}
    \label{tab:flip-a2r}
    \small
    \begin{tabular}{p{7.8cm}lp{2.8cm}}
        \toprule
        Condition & A2 & Claude's label \tabularnewline
        \midrule
        No edit (stored) & 1884 & adoption \tabularnewline
        Targeted edit, pretrained SAE, \textsc{decrease} & 1927 & \textbf{retention} (flip) \tabularnewline
        Random edits 1--3, same SAE and arm & 1884 & adoption \tabularnewline
        \bottomrule
    \end{tabular}
\end{table}

Claude's reason for the targeted answer: ``A1 states 1927, C1 states 1884
(different positions), and A2 maintains 1927, matching A1's position despite
the critic's feedback.''

\begin{table}[htbp]
    \centering
    \caption{A Retention-to-Adoption flip (MuSiQue, two-hop). Question: ``When
    did the performer of \emph{Out from Under} release her first album?'' A1:
    1999. C1: 2000. All edits have the same size, $0.2\lVert x\rVert$.}
    \label{tab:flip-r2a}
    \small
    \begin{tabular}{p{7.8cm}lp{2.8cm}}
        \toprule
        Condition & A2 & Claude's label \tabularnewline
        \midrule
        No edit (stored) & 1999 & retention \tabularnewline
        Targeted edit, decoder-adapted SAE, \textsc{both} & 2000 & \textbf{adoption} (flip) \tabularnewline
        Random edits 1--3, same SAE and arm & 1999 & retention \tabularnewline
        \bottomrule
    \end{tabular}
\end{table}

Claude's reason for the targeted answer: ``A1 states 1999, C1 states 2000,
and A2 changes from A1's position to match C1's position of 2000.''

In both episodes the edited answer is the other candidate, written the same
way, not degraded text. A broken answer looks different. Of the 60 edited
answers to the Marnhill question across all conditions, two were labeled
Out. Both came from the same random draw of the raw-activation control, under
the \textsc{increase} and \textsc{both} arms, and both were simply ``8''.
Claude's reason: ``A2 states `8' which does not match A1's position of `1927'
or C1's position of `1884'. The response appears to be neither a year nor a
semantically equivalent value to either candidate position.''
